% !PRINTED 583
\weberpage{583}

%<<CONT
worin
\[
H_{l,v'} + \Omega_{h,h'} = \dfrac{2(lm+hn)K + 2(lm+h'n)iK'}{mn}, % [source sic: the second coefficient prints lm+h'n, although the definition gives l'm+h'n]
\]
und folglich sind die Faktoren des Produktes (11):
\[
\sqrt{\varkappa}\ \mathrm{sn}\,(H_{l,v'} + \Omega_{h,h'})
\]
ganze Zahlen [nach (6)]. Vertauscht man in (11) wieder $m$ und $n$, so folgt:
\weq{12}{\dfrac{\sqrt{\varkappa}\ \mathrm{sn}\,n\,\Omega_{h,h'}}{\sqrt{\varkappa}\ \mathrm{sn}\,\Omega_{h,h'}} \text{ ist eine ganze Zahl.}}
Dies gilt für beliebige ungerade $n$, also auch, wenn $n$ relativ prim zu $m$ ist, für $n^{-1}$ (mod $m$) und es ist also auch
\[
\dfrac{\sqrt{\varkappa}\ \mathrm{sn}\,n^{-1}\,\Omega_{h,h'}}{\sqrt{\varkappa}\ \mathrm{sn}\,\Omega_{h,h'}} \text{ eine ganze Zahl.}
\]
Ersetzt man hierin $h,h'$ durch $nh, nh'$, so folgt, da\ss{} auch
\[
\dfrac{\sqrt{\varkappa}\ \mathrm{sn}\,\Omega_{h,h'}}{\sqrt{\varkappa}\ \mathrm{sn}\,n\,\Omega_{h,h'}} \text{ eine ganze Zahl ist,}
\]
also ist der Quotient (12) eine \gesperrt{Einheit}, oder anders ausgedrückt:

2. \gesperrt{Durchläuft $n$ eine Reihe ungerader, zu $m$ teilerfremder Zahlen, so sind die ganzen Zahlen:}
\weq{13}{\sqrt{\varkappa}\ \mathrm{sn}\,n\,\Omega_{h,h'}}
\gesperrt{miteinander assoziiert.}

\begin{center}
\textbf{\S\,157.\qquad \gesperrt{Komplexe Multiplikatoren.}}
\end{center}

Es genüge $\omega$ wie im vorigen Paragraphen der quadratischen Gleichung
\weq{1}{a\omega^2 + b\omega + c = 0}
mit negativer Stammdiskriminante
\[
\varDelta = b^2 - 4ac,
\]
also sei:
\weq{2}{\omega = \dfrac{-b + \sqrt{\varDelta}}{2a}.}
Wir legen unseren Betrachtungen wie im vorigen Paragraphen die Funktion
\weq{3}{\sqrt{\varkappa}\ \mathrm{sn}\,v = \dfrac{\vartheta_{11}(u)}{\vartheta_{01}(u)} = \boldsymbol{s}(u)}
%CONT>>
